\section{Comparison with the arithmetic character question}

The abelian quotients constructed here measure the effect of omitting
the pentagon. The arithmetic question concerns a different restriction:
a character on full \(\GT\) that vanishes on its Galois subgroup.
The companion paper \cite{ArithmeticCompanion} studies that restriction
using the same cyclic-cover homology, together with CM reciprocity.
We state its conclusions with their precise scope.

Write \(G=\GT\), \(H=G_{\mathbb Q}\), and
\[
 N=\ker(\lambda,\rho_2),\qquad K=N\cap H,\qquad
 B=N/\overline{[G,N]K}.
\]
Here \(\rho_2\) is the Kummer cocycle of two in the normalization
of \cite{H17}. The decomposition \(G=NH\) gives
\[
 G^{\rm ab}\simeq\widehat{\mathbb Z}^{\times}\times B.
\]
In particular, a nonzero continuous quadratic character of \(B\)
would separate \(G\) from \(H\).
The invariant quadratic restriction
\[
 H^1(N,\mathbb F_2)^G\longrightarrow H^1(K,\mathbb F_2)^H
\]
is surjective and has an explicit linear section
\cite[Theorems~1.1 and~1.6]{ArithmeticCompanion}.
Its kernel is \(\operatorname{Hom}_{\rm cts}(B,\mathbb F_2)\),
whose vanishing remains unresolved.

The additional CM comparisons give integral arithmetic normalizations
of this same relative obstruction
\cite[Theorem~1.14]{ArithmeticCompanion}.
For \(y^p=x(1-x)\), a local reflex-norm operator determines a finite
deficit in a dyadic norm-one torus. Its determinant is governed by
the relative cyclotomic class number, but the complete arithmetic
image also depends on square classes of explicit real units.
Frobenius elements fill the deficits at \(p=29\) and \(113\).
At \(p=163\), negative-square identities prove that the arithmetic
image still has index four. Simultaneous arithmetic matching is used
only for the coordinates whose entire images are known to match.
For any collection of prime degrees, the full and arithmetic scalar
images are nevertheless exact products of their individual images.
The entire conjugation obstruction for intrinsic relative quadratic
characters vanishes, including for the infinite collection. Such a
character extends to full \(G\), zero on \(H\), precisely when its
relative central extension class is zero
\cite[Theorem~1.16]{ArithmeticCompanion}.

These results determine the arithmetic target and its lifted part
without determining the remaining relative kernel. The characters
constructed in the present paper vanish on all of \(\GT\), so their
restrictions to that kernel are zero. They prove strictness of the
gentle inclusion and provide explicit abelian quotients of its
normal quotient; an arithmetic-trivial character nonzero on full
\(\GT\) would require an additional argument.
